\title{DeepSeekMath: Pushing the Limits of Mathematical Reasoning in Open Language Models}

\begin{abstract}
Mathematical reasoning remains a notable obstacle for language models due to its inherently intricate and highly structured characteristics. In this work, we present DeepSeekMath~7B, developed by further pre-training DeepSeek-Coder-Base-v1.5 7B with an additional 120B math-focused tokens derived from Common Crawl, alongside diverse natural language and programming code sources. DeepSeekMath~7B achieves a remarkable 51.7\% accuracy on the competition-grade MATH benchmark without the assistance of external toolkits or ensembling strategies, approaching the performance of leading models like Gemini-Ultra and GPT-4. Utilizing a self-consistency protocol with 64 generated samples, DeepSeekMath~7B attains a 60.9\% score on MATH. The advances in DeepSeekMath's mathematical reasoning can be traced to two principal components: Firstly, we exploit the extensive potential of freely available online data through a carefully designed data selection framework. Secondly, we propose Group Relative Policy Optimization (GRPO)—an adaptation of Proximal Policy Optimization (PPO)—which bolsters mathematical reasoning competence while simultaneously improving PPO's memory efficiency.
\end{abstract}

\section{Introduction}

The advent of large language models (LLM) has transformed how artificial intelligence approaches mathematical reasoning, catalyzing substantial progress in both the quantitative reasoning benchmark \citep{MATH} and the geometry reasoning benchmark \citep{trinh2024solving}. Additionally, these systems have become valuable tools for aiding humans in tackling sophisticated mathematical tasks \citep{tao}. Despite these achievements, state-of-the-art models, including GPT-4 \citep{gpt4} and Gemini-Ultra \citep{gemini}, remain inaccessible to the public, with available open-source alternatives displaying significantly inferior performance.

This paper presents DeepSeekMath, a specialized language model tailored for mathematical tasks, which demonstrates a marked improvement over existing open-source models and nearly matches the performance of GPT-4 in standardized academic evaluations. Central to this advancement is the creation of the DeepSeekMath Corpus—a comprehensive pre-training dataset consisting of 120B mathematics-related tokens. The construction of this dataset begins with extraction from Common Crawl (CC) via a classifier based on fastText \citep{joulin2016fasttext}. Initially, the classifier is trained using examples from OpenWebMath \citep{openwebmath} as the positive class and a broad range of unrelated web pages as negative samples. This preliminary classifier is then employed to retrieve further positive samples from the CC, which are subsequently validated and refined by human annotators. Incorporating these newly annotated examples, the classifier undergoes additional training to enhance its accuracy. Empirical analysis confirms the exceptional quality of the resultant corpus, as evidenced by DeepSeekMath-Base 7B attaining 64.2\% accuracy on GSM8K \citep{gsm8k} and 36.2\% on the advanced MATH benchmark \citep{MATH}, outperforming models like Minerva 540B \citep{minerva}. Notably, the DeepSeekMath Corpus is designed to support multiple languages, which also leads to observed gains on Chinese mathematical datasets \citep{wei2023cmath,agieval}. We contend that our methodology for constructing and refining mathematical corpora may serve as a valuable foundation for subsequent research, with ample opportunity for future enhancements.

DeepSeekMath-Base is built upon the DeepSeek-Coder-Base-v1.5 7B model \citep{deepseek-coder}, as our findings suggest that initializing from a code-pretrained model yields superior results compared to leveraging a generic large language model. Additionally, we note that training with mathematical data not only boosts performance in math-specific domains but also enhances results on broader evaluation benchmarks such as MMLU \citep{mmlu} and BBH \citep{bbh}. This observation implies an overall improvement in both mathematical proficiency and general reasoning ability.

Following pre-training, we conduct instruction tuning for DeepSeekMath-Base using mathematical datasets that encompass chain-of-thought reasoning \citep{cot}, program-of-thought approaches \citep{pot,pal}, and tool-augmented reasoning methods \citep{tora}. The output model, DeepSeekMath-Instruct 7B, demonstrates superiority over all existing 7B models and matches the performance of some of the leading 70B instruction-tuned open-source models.

We also present Group Relative Policy Optimization (GRPO), a reinforcement learning algorithm derived from Proximal Policy Optimization (PPO) \citep{schulman2017proximal}. Unlike traditional PPO, GRPO eliminates the need for a critic network, opting instead to calculate baselines via group-level scoring, leading to significant reductions in computational demands. Leveraging only a subset of English instruction data during reinforcement learning, GRPO markedly advances performance over DeepSeekMath-Instruct—evidenced by notable gains in both domain-specific (GSM8K: 82.9\% $\rightarrow$ 88.2\%, MATH: 46.8\% $\rightarrow$ 51.7\%) and transfer tasks (such as CMATH: 84.6\% $\rightarrow$ 88.8\%). 

Moreover, we outline a unified framework that encompasses a variety of related methods, including Rejection Sampling Fine-Tuning (RFT) \citep{yuan2023scaling}, Direct Preference Optimization (DPO) \citep{dpo}, PPO, and our proposed GRPO. Within this framework, we interpret these approaches as variants of direct or simplified reinforcement learning strategies. Through comprehensive experimentation—comparing online versus offline paradigms, outcome versus process-based supervision, and single versus iterative reinforcement learning—we delve into the fundamental factors shaping this paradigm. Finally, we elucidate how reinforcement learning enhances instruction-tuned model performance and discuss prospective avenues for advancing RL efficacy guided by this unified perspective.

\subsection{Contributions}
This work presents advances in large-scale math pre-training and investigates the application and effects of reinforcement learning.

\textbf{Math Pre-Training at Scale}

\begin{itemize}[topsep=0pt]
    \item Our study demonstrates that the Common Crawl dataset, which is publicly available, holds significant mathematical content. Through a rigorously developed data selection process, we assemble the DeepSeekMath Corpus—an extensive and high-quality dataset encompassing 120B tokens drawn from math-focused web content. This corpus is approximately seven times larger than the dataset utilized by Minerva \citep{minerva}, and nine times larger than the dataset in OpenWebMath \citep{openwebmath}.

    \item The DeepSeekMath-Base 7B model, pre-trained on this dataset, reaches performance levels comparable to those of Minerva 540B \citep{minerva}. This outcome suggests that model size alone does not dictate success in mathematical reasoning tasks; strong results can be attained by relatively smaller models trained on superior-quality data.

    \item We report insights from our math training trials. Incorporating code training prior to math-focused pre-training leads to enhanced mathematical problem-solving, both in the presence and absence of tool assistance. These observations contribute evidence toward resolving the ongoing debate regarding whether code training enhances reasoning, at least in mathematical domains, supporting the idea that it does.

    \item In contrast to common practice, incorporating arXiv papers in the training process does not yield measurable gains across any of the mathematical benchmarks evaluated in this research.
\end{itemize}

\textbf{Exploration and Analysis of Reinforcement Learning}

\begin{itemize}[topsep=0pt]
    \item We present Group Relative Policy Optimization (GRPO), a novel reinforcement learning algorithm that prioritizes computational efficiency and effectiveness. Rather than employing a critic model, GRPO utilizes group score-based baseline estimation, thereby cutting down on the resource demands commonly associated with Proximal Policy Optimization (PPO).
    \item We provide empirical evidence that GRPO significantly improves the outcomes of our instruction-tuned DeepSeekMath-Instruct model, relying exclusively on instruction-tuning datasets. Additionally, the reinforcement learning process with GRPO leads to improved out-of-domain performance.
    \item Our work offers a comprehensive framework that unifies the understanding of approaches such as RFT, DPO, PPO, and GRPO. We perform an in-depth experimental investigation covering scenarios like online versus offline training, outcome-based versus process-based supervision, and single-step versus iterative reinforcement learning, shedding light on the key factors underlying this framework.
    \item Leveraging our unified framework, we analyze the mechanisms that underpin the efficacy of reinforcement learning and outline multiple promising avenues for further enhancing the reinforcement learning of LLMs.
\end{itemize}

\subsection{Summary of Evaluations and Metrics}

\begin{itemize}[topsep=0pt]
    \item \textbf{English and Chinese Mathematical Reasoning}: Our evaluation encompasses a wide range of English and Chinese benchmarks, assessing mathematical reasoning abilities across problems ranging from elementary to university level. The English benchmarks involve GSM8K \citep{gsm8k}, MATH \citep{MATH}, SAT \citep{llemma}, OCW Courses \citep{minerva}, and MMLU-STEM \citep{mmlu}. Chinese benchmarks comprise MGSM-zh \citep{mgsm}, CMATH \citep{wei2023cmath}, Gaokao-MathCloze \citep{agieval}, and Gaokao-MathQA \citep{agieval}. The assessment considers both the models’ competence in generating fully self-contained text-based solutions without external tool assistance and their ability to solve problems using Python.

    In the context of English benchmarks, DeepSeekMath-Base exhibits comparable performance to the proprietary Minerva 540B \citep{minerva}, and substantially outperforms all open-source foundation models (such as Mistral 7B \citep{mistral} and Llemma-34B \citep{llemma}), irrespective of prior math pre-training, with especially large gains in many cases. Remarkably, DeepSeekMath-Base attains leading results on Chinese datasets. This may be attributed to our approach diverging from earlier works \citep{minerva,llemma} by incorporating high-quality non-English math data into pre-training. When augmented with mathematical instruction tuning and reinforcement learning, DeepSeekMath-Instruct and DeepSeekMath-RL both achieve robust performance, surpassing the 50\% accuracy threshold for the first time on the MATH competition benchmark in the open-source landscape.
\end{itemize}

\item \textbf{Formal Mathematics}: For the evaluation of DeepSeekMath-Base, we adopt the informal-to-formal theorem proving task as specified in \citep{dsp_proof}, utilizing miniF2F \citep{minif2f} and selecting Isabelle \citep{isabelle} as the proof assistant. DeepSeekMath-Base exhibits notable performance in few-shot autoformalization settings.
\item \textbf{Natural Language Understanding, Reasoning, and Code}: To establish a broad assessment of the model’s capabilities in general comprehension, logical reasoning, and programming, we benchmark DeepSeekMath-Base on Massive Multitask Language Understanding (MMLU) \citep{mmlu}, comprising 57 diverse multiple-choice subjects; BIG-Bench Hard (BBH) \citep{bbh}, a set of 23 intricate tasks largely dependent on multi-step reasoning; as well as HumanEval \citep{codex} and MBPP \citep{mbpp}, which serve as standard datasets for evaluating code generation models. Pre-training on mathematical corpora contributes positively to both language understanding and reasoning proficiency.

\section{Math Pre-Training}

\subsection{Data Collection and Decontamination} This section details our methodology for assembling the DeepSeekMath Corpus using Common Crawl data. As illustrated in Figure \ref{fig:data_collect}, our approach employs a stepwise pipeline to reliably compile an extensive collection of mathematical documents from Common Crawl, beginning with an initial, high-quality seed set—such as a well-curated mathematics dataset. Notably, this framework is extensible to various domains, including but not limited to programming.

Initially, we designate OpenWebMath \citep{openwebmath}, which provides high-caliber mathematical web content, as the seed corpus. Leveraging this set, we train a fastText model \citep{joulin2016fasttext} to retrieve additional web pages that resemble the OpenWebMath dataset in mathematical content. Specifically, 500,000 examples are randomly chosen from the seed set as positive instances, with another 500,000 web pages from Common Crawl serving as negative samples. The training process uses an open-source implementation\footnote{\url{https://fasttext.cc}}, employing the following hyperparameters: vector dimensionality of 256, a learning rate of 0.1, a maximum word n-gram length of 3, minimum occurrence count of 3 for words, and three training epochs. To minimize data redundancy in the Common Crawl source, we apply both URL-level deduplication and near-duplicate filtering, which yields a pool of 40B unique HTML pages. Subsequently, mathematical content is retrieved from this deduplicated set using the trained fastText model. To ensure data quality, retrieved web pages are scored by the fastText classifier, after which only those with the highest rankings are retained. The selected dataset size is calibrated through pre-training trials with the top 40B, 80B, 120B, and 160B tokens. For the initial phase, we opt to retain the top 40B tokens.

Following the initial data collection phase, a substantial number of mathematical web pages remain unacquired, largely due to the limited diversity within the set of positive samples used to train the fastText model. To enhance the model, we seek out further mathematical web domains to expand the seed corpus. Our procedure begins with partitioning the entire Common Crawl dataset into distinct domains, with each domain defined as the collection of web pages sharing a common base URL. For each domain, we compute the proportion of web pages captured during the first collection round. If more than 10\% of a domain’s pages are gathered, the domain is flagged as being math-related (for instance, \url{mathoverflow.net}).

We then carry out manual annotation to identify URLs within these domains that specifically host mathematical content (for example, \url{mathoverflow.net/questions}). Pages corresponding to these URLs that have not yet been collected are incorporated into the seed corpus. This strategy increases the volume of positive samples, enabling the training of a more effective fastText model, which in turn yields higher recall of mathematical data in subsequent iterations. After four rounds of data collection, the total comprises 35.5M mathematical web pages and sums to 120B tokens. Notably, by the fourth collection phase, approximately 98\% of the dataset was already assembled by the third iteration, prompting us to halt further data acquisition.

To prevent contamination from benchmarks, we adopt the approach described in \cite{deepseek-coder} to exclude web pages that include content from mathematical evaluation datasets in English (such as GSM8K \citep{gsm8k} and MATH \citep{MATH}) and Chinese (including CMATH \citep{wei2023cmath} and AGIEval \citep{agieval}). The filtering procedure operates as follows: any segment of text containing an exact match for any 10-gram string from an evaluation benchmark is discarded from our mathematical training data. For benchmark entries that are fewer than 10 but at least 3 tokens long, precise string matching is used to identify and remove contaminated web content.

\subsection{Validating the Quality of the DeepSeekMath~Corpus}

We conduct pre-training studies to compare the DeepSeekMath~Corpus to recently developed mathematical datasets:

\begin{itemize}[topsep=0pt]
    \item \textbf{MathPile} \citep{mathpile}: an 8.9B-token collection drawn from multiple sources—including textbooks, Wikipedia, ProofWiki, CommonCrawl, StackExchange, and arXiv—with arXiv constituting over 85\% of the total;
    \item \textbf{OpenWebMath} \citep{openwebmath}: a corpus of 13.6B tokens extracted from CommonCrawl and curated for mathematical relevance;
    \item \textbf{Proof-Pile-2} \citep{llemma}: a dataset comprising OpenWebMath, 10.3B tokens of mathematical code from AlgebraicStack, and 28.0B tokens from arXiv papers. For experiments utilizing Proof-Pile-2, we employ the arXiv:Web:Code ratio of 2:4:1 as specified in \cite{llemma}.
\end{itemize}

\subsubsection{Training Setting}
\label{sec:quality-policy}
We adapt a general-purpose language model, equipped with 1.3B parameters and architecturally consistent with the DeepSeek LLMs \citep{deepseek-llm}—referred to here as DeepSeek-LLM 1.3B—for specialized mathematical training. Separate instances of this model are fine-tuned on distinct mathematical datasets, each over a span of 150B tokens. All model training is conducted using HAI-LLM \citep{haillm}, an efficient and lightweight training infrastructure. Mirroring the training strategy of DeepSeek LLMs, we utilize the AdamW optimizer \citep{adamW}, setting $\beta_1=0.9$, $\beta_2=0.95$, and $\mathrm{weight\_decay}=0.1$. The learning rate follows a multi-stage decay schedule: after 2,000 warmup steps, it attains its maximum, then declines to 31.6\% at 80\% of the training duration, and subsequently drops to 10.0\% of its peak at 90\% completion. The peak learning rate is capped at 5.3e-4, with batch processing set to 4M tokens per step and a context window of 4K tokens.

\subsubsection{Evaluation Results}

\textbf{The DeepSeekMath~Corpus demonstrates exceptional quality, encompasses multilingual mathematical material, and surpasses other corpora in scale.}

\begin{itemize}[topsep=0pt]
    \item \textbf{High-quality}: 
    We assess downstream capabilities using eight mathematical evaluation sets, employing few-shot chain-of-thought prompting \cite{cot}. Table \ref{tab:corpora_comparison} highlights a pronounced advantage for models trained with the DeepSeekMath~Corpus. As illustrated in Figure \ref{fig:corpora_comparisons}, DeepSeekMath-trained models outperform those trained on Proof-Pile-2 when evaluated at 50B tokens (corresponding to a single epoch on Proof-Pile-2), signifying that DeepSeekMath delivers a higher average data quality.
    \item \textbf{Multilingual}: 
    DeepSeekMath~Corpus contains mathematical texts in various languages, with English and Chinese comprising the largest portions. Training with this dataset noticeably boosts mathematical reasoning accuracy in both languages, as shown in Table \ref{tab:corpora_comparison}. By contrast, prior datasets—predominantly focused on English—provide minimal benefit or may even impair Chinese mathematical reasoning abilities.
    \item \textbf{Large-scale}: 
    In terms of dataset volume, DeepSeekMath~Corpus considerably exceeds current mathematical corpora. Figure \ref{fig:corpora_comparisons} depicts DeepSeek-LLM 1.3B's rapid and persistent performance growth when fine-tuned on DeepSeekMath, while smaller comparison corpora are reused extensively during training, resulting in stagnating performance.
\end{itemize}

\subsection{Training and Evaluating DeepSeekMath-Base 7B}

This section presents DeepSeekMath-Base 7B, a foundational model demonstrating robust reasoning capabilities, particularly within mathematical domains. The model's initial weights are derived from DeepSeek-Coder-Base-v1.5 7B \citep{deepseek-coder}, after which it undergoes training over a corpus totaling 500B tokens. The dataset composition is as follows: DeepSeekMath Corpus constitutes 56\%, AlgebraicStack accounts for 4\%, arXiv contributes 10\%, Github code comprises 20\%, and the last 10\% is sourced from natural language content in both English and Chinese from Common Crawl. The majority of the training configuration aligns with Section \ref{sec:quality-policy}, with the exceptions that the peak learning rate is set at 4.2e-4, and the batch size is increased to 10M tokens.

A thorough evaluation is undertaken to analyze DeepSeekMath-Base 7B's mathematical proficiency. We assess its competence in independently generating comprehensive mathematical solutions, solving problems via tool use, and performing formal theorem verification. Additionally, we report on its broader capabilities, such as natural language understanding, reasoning, and programming performance, to provide an overall characterization of the base model.

\paragraph{Mathematical Problem Solving with Step-by-Step Reasoning}

To assess DeepSeekMath-Base's ability to solve mathematical problems, we employ few-shot chain-of-thought prompting \citep{cot} across eight different English and Chinese benchmarks. These cover a wide spectrum of mathematical problem types, including quantitative reasoning tasks such as GSM8K \citep{gsm8k}, MATH \citep{MATH}, and CMATH \citep{wei2023cmath}, as well as multiple-choice formats like MMLU-STEM \citep{mmlu} and Gaokao-MathQA \citep{agieval}. The benchmarks span various mathematical domains, from elementary-level questions up to challenges encountered in undergraduate studies.

As detailed in Table \ref{tab:base_math_cot}, among open-source base models—including the prevalent Mistral 7B \citep{mistral} and the more recent Llemma 34B \citep{llemma}, which was specifically trained on the Proof-Pile-2 dataset \citep{llemma} for mathematics—DeepSeekMath-Base 7B achieves the highest results across all eight evaluated benchmarks. On the challenging MATH dataset at the competition level, DeepSeekMath-Base 7B exceeds the next-best open-source base models by more than 10 percentage points in absolute terms, and even surpasses the much larger, closed-source Minerva 540B \citep{minerva}, a model that is 77 times bigger, is based on PaLM \citep{palm}, and has received additional training on mathematical materials.

\paragraph{Mathematical Problem Solving with Tool Use} For the assessment of program-assisted mathematical reasoning, we apply few-shot program-of-thought prompting \citep{pot,pal} to both GSM8K and MATH datasets. The evaluation setup directs the models to develop Python solutions for each given problem, with the flexibility to employ libraries such as \textit{math} and \textit{sympy} for handling complex calculations. The correctness of a solution is determined by executing the produced code and inspecting its output. As indicated in Table \ref{tab:base_math_pot_proof}, DeepSeekMath-Base 7B establishes a new open-source benchmark by outperforming Llemma 34B, the previous state-of-the-art model in this category.

\paragraph{Formal Mathematics}

The automation of formal proofs plays an essential role in improving both the reliability and efficiency of mathematical reasoning, which has recently drawn substantial interest. We benchmark DeepSeekMath-Base 7B using the informal-to-formal proof generation task described in \citep{dsp_proof}. Here, the model is tasked with generating formal proofs by leveraging an informal problem statement, its formal rephrasing, and an informal proof as input. On the miniF2F \citep{minif2f} benchmark—a testbed for Olympiad-level mathematics formalization—we prompt the model to generate proofs in Isabelle using a few-shot learning approach. Following the methodology in \cite{dsp_proof}, generated proof outlines are further refined using the automated theorem prover Sledgehammer \citep{sledgehammer} to complete unresolved components. Table \ref{tab:base_math_pot_proof} shows that DeepSeekMath-Base 7B delivers robust results in automated formal proof construction.

\paragraph{Natural Language Understanding, Reasoning, and Code}

We examine natural language comprehension on MMLU \citep{mmlu}, reasoning abilities using BBH \citep{bbh}, and programming capabilities with HumanEval \citep{codex} and MBPP \citep{mbpp}. Table \ref{tab:understanding_reasoning_code} reveals that DeepSeekMath-Base 7B considerably improves over its predecessor, DeepSeek-Coder-Base-v1.5 \citep{deepseek-coder}, on both MMLU and BBH tasks, underscoring the beneficial effect of math-focused training on broader language and reasoning competencies. Moreover, by retaining code token coverage in ongoing training, DeepSeekMath-Base 7B sustains the coding proficiency of DeepSeek-Coder-Base-v1.5 on the HumanEval and MBPP benchmarks. In summary, DeepSeekMath-Base 7B provides substantial performance gains compared to the general-purpose Mistral 7B model \citep{mistral} across all evaluated reasoning and coding tasks.

\section{Supervised Fine-Tuning}

\subsection{SFT Data Curation}

We develop a dataset for mathematical instruction-tuning that spans both English and Chinese, drawing from numerous mathematical domains and encompassing a range of difficulties. Each problem in the dataset is accompanied by a corresponding solution, formatted as either chain-of-thought (CoT) \citep{cot}, program-of-thought (PoT) \citep{pot,pal}, or with integrated tool-based reasoning \citep{tora}. In total, the dataset consists of 776,000 training instances.

\begin{itemize}[topsep=0pt]
    \item \textbf{English mathematical datasets}: We provide tool-integrated solution annotations for GSM8K and MATH, and additionally leverage portions of MathInstruct \citep{MathInstruct} and the Lila-OOD training corpus \citep{lila}, which offer problems solved via CoT or PoT. The English portion of the dataset represents a broad selection of mathematical fields, such as algebra, calculus, geometry, probability, and number theory.
    \item \textbf{Chinese mathematical datasets}: We compile Chinese-language math questions from the K-12 curriculum, encompassing 76 distinct subtopics like linear equations, and supply both CoT and tool-integrated solution annotations.
\end{itemize}

\subsection{Training and Evaluating DeepSeekMath-Instruct 7B}

Here, we detail the mathematical instruction-tuning process for DeepSeekMath-Instruct 7B, which is further trained on DeepSeekMath-Base. During fine-tuning, training samples are concatenated at random up to a total context window of 4K tokens. The model is optimized for 500 training steps using a batch size of 256 and a constant learning rate set at 5e-5.

To assess mathematical capabilities, we evaluate models on four quantitative reasoning benchmarks in both English and Chinese, with and without access to tool integration. Comparative benchmarking is conducted against the most advanced models available at the time:

\begin{itemize}[topsep=0pt]
    \item \textbf{Closed-source models} assessed are: (1) the GPT series, where GPT-4 \citep{gpt4} and the GPT-4 Code Interpreter~\footnote{\url{https://openai.com/blog/chatgpt-plugins\#\#code-interpreter}} represent top performance, (2) Gemini Ultra and Pro \citep{gemini}, (3) Inflection-2 \citep{inflection-2}, (4) Grok-1~\footnote{\url{https://x.ai/model-card}}, as well as newer releases from Chinese providers, including (5) Baichuan-3~\footnote{\url{https://www.baichuan-ai.com}}, and (6) the latest GLM-4~\footnote{\url{https://open.bigmodel.cn/dev/api\#glm-4}} from the GLM model family \citep{glm}.
\end{itemize}

The aforementioned models are primarily designed for broad applications, with most undergoing multiple alignment processes.

\item \textbf{Open-source models} span general-purpose architectures such as (1) DeepSeek-LLM-Chat 67B \citep{deepseek-llm}, (2) Qwen 72B \citep{qwen}, (3) SeaLLM-v2 7B \citep{seallm}, and (4) ChatGLM3 6B \citep{chatglm3}. Additionally, several models are tailored for mathematical tasks: (5) InternLM2-Math 20B~\footnote{\url{https://github.com/InternLM/InternLM-Math}} extends InternLM2 with additional mathematics training and instruction fine-tuning; (6) Math-Shepherd-Mistral 7B leverages PPO training \citep{schulman2017proximal} on top of Mistral 7B \citep{mistral} and incorporates a process-supervised reward system; (7) the WizardMath series \citep{wizardmath} advances mathematical reasoning capabilities in both Mistral 7B and Llama-2 70B \citep{llama2} by combining evolve-instruct (a variant of instruction tuning employing AI-generated prompts) with PPO training, primarily utilizing GSM8K and MATH datasets; (8) MetaMath 70B \citep{metamath} consists of Llama-2 70B further trained on an expanded GSM8K and MATH; (9) ToRA 34B \cite{tora} adapts CodeLlama 34B for integrated tool-assisted mathematical reasoning via fine-tuning; and (10) MAmmoTH 70B \citep{MathInstruct} involves instruction-tuning Llama-2 70B using MathInstruct.

Table \ref{tab:sft_rl_math} presents empirical results under the constraint of disallowing tool usage, where DeepSeekMath-Instruct 7B exhibits notably robust sequential reasoning abilities. Especially on the challenging, competition-style MATH benchmark, this model exceeds the accuracy of all open-source competitors as well as most commercial models (such as Inflection-2 and Gemini Pro) by margins of at least 9\% absolute. This advantage persists even relative to considerably larger systems (e.g., Qwen 72B) or models specifically refined for mathematics through reinforcement learning (for instance, WizardMath-v1.1 7B). While DeepSeekMath-Instruct is on par with leading Chinese proprietary offerings such as GLM-4 and Baichuan-3 for the MATH dataset, it still lags behind state-of-the-art solutions like GPT-4 and Gemini Ultra.

In the alternative evaluation scenario that permits models to utilize both natural language inference and code-based tools for solving tasks, DeepSeekMath-Instruct 7B achieves nearly 60\% accuracy on MATH, thereby outperforming all currently available open-source alternatives. For other benchmarks, it matches the performance of DeepSeek-LLM-Chat 67B—the prior leading model, which is tenfold larger in size.

\section{Reinforcement Learning}

\subsection{Group Relative Policy Optimization} Previous research has demonstrated that reinforcement learning (RL) is effective at further refining the mathematical reasoning capability of large language models following Supervised Fine-Tuning (SFT) \citep{wang2023math,wizardmath}. Here, we describe our novel RL approach: Group Relative Policy Optimization (GRPO), designed for efficiency and effectiveness.

\subsubsection{From PPO to GRPO} Proximal Policy Optimization (PPO) \citep{schulman2017proximal}, a widely adopted actor-critic RL technique, serves as the standard method for RL-based tuning in LLMs \citep{ouyang2022training}. PPO updates model policies by maximizing the surrogate objective given by:

\begin{equation}
\footnotesize
    \mathcal{J}_{PPO}(\theta) = \mathbb{E}{[q \sim P(Q), o \sim \pi_{\theta_{old}}(O|q)]} \frac{1}{|o|} \sum_{t=1}^{|o|} \min \left[ \frac{\pi_\theta(o_{t} | q, o_{<t})}{\pi_{\theta_{old}}(o_{t} | q, o_{<t})} A_{t}, \text{clip} \left( \frac{\pi_\theta(o_{t} | q, o_{<t})}{\pi_{\theta_{old}}(o_{t} | q, o_{<t})}, 1 - \varepsilon, 1 +

Here, $\pi_{\theta}$ and $\pi_{\theta_{old}}$ denote the current and previous policy models, respectively, while $q$ and $o$ represent questions and outputs, which are sampled from the question dataset and according to the previous policy $\pi_{\theta_{old}}$. The hyper-parameter $\varepsilon$ is associated with the clipping mechanism in PPO, introduced to enhance training stability. The term $A_t$ stands for the advantage estimate, calculated using Generalized Advantage Estimation (GAE) \citep{gae}, which leverages rewards $\{r_{\ge t}\}$ and a value function $V_{\psi}$ that is learned. Consequently, PPO requires the simultaneous training of a value function alongside the policy model. To counteract reward model over-optimization, the common practice involves integrating a per-token KL penalty from a reference model into the reward at every token step \citep{ouyang2022training}, such that:

\begin{equation}
    r_{t} = r_\varphi(q, o_{\le t}) - \beta \log\frac{\pi_{\theta}(o_{t}|q, o_{<t})}{\pi_{ref}(o_{t}|q, o_{<t})},
\label{eq:PPO-reward}
\end{equation}

Here, $r_\varphi$ is the reward model; $\pi_{ref}$ refers to a reference policy, commonly chosen as the original SFT model; and $\beta$ regulates the KL penalty term.

Since the value function typically implemented in PPO is itself a model on par with the policy model in size, this setup imposes considerable demands on memory and computation. Furthermore, during RL optimization, the value function acts as a baseline to reduce the variance of the advantage estimates. In the case of large language models (LLMs), the reward model usually provides a reward only for the final token, which poses a challenge when training a value function that accurately predicts values at every token step. To overcome these limitations, as depicted in Figure \ref{fig:grpo}, we introduce Group Relative Policy Optimization (GRPO). This approach eliminates the need for learning an explicit value function (as is necessary in PPO) by using, instead, the mean reward of several sampled outputs—generated for the same question—as a baseline. Specifically, for each question $q$, GRPO draws a set of outputs $\{o_1, o_2, \cdots, o_G\}$ from the old policy $\pi_{\theta_{old}}$ and then updates the policy by maximizing the following objective:

\begin{equation}
\footnotesize
\begin{split}
    \mathcal{J}_{GRPO}(\theta) &= \mathbb{E}{[q \sim P(Q), \{o_i\}_{i=1}^G \sim \pi_{\theta_{old}}(O|q)]}  \\
    & \frac{1}{G}\sum_{i=1}^G\frac{1}{|o_i|} \sum_{t=1}^{|o_i|} \left\{ \min \left[ \frac{\pi_\theta(o_{i,t} | q, o_{i,<t})}{\pi_{\theta_{old}}(o_{i,t} | q, o_{i,<t})} \hat{A}_{i,t}, \text{clip} \left( \frac{\pi_\theta(o_{i,t} | q, o_{i,<t})}{\pi_{\theta_{old}}(o_{i,t} | q, o_{i,<t})}, 1 - \varepsilon, 1 + \varepsilon \right)  \hat{A}_{i,t} \right] - \beta \mathbb{D}_{KL}\left[\pi_{\theta} || \pi_{ref}\right]\right\} ,
\end{split}
\label{eq:GRPO-obj}
\end{equation}

where $\varepsilon$ and $\beta$ serve as hyper-parameters, and $\hat{A}_{i,t}$ is the advantage computed using only the relative rewards among the group of sampled outputs, as will be further explained in subsequent sections. GRPO's method of calculating advantages by comparing samples within a group is well-suited to the comparative foundation of most reward models, as these models are typically trained on pairwise comparisons between outputs to the same question. Additionally, GRPO departs from reward-level KL penalization by incorporating the KL divergence between the updated policy and the reference model directly in the objective function, simplifying the computation of $\hat{A}_{i,t}$. Importantly, while Equation~(\ref{eq:PPO-reward}) applies a per-token KL penalty, we instead

\begin{algorithm}[t]
  \small
  \caption{Iterative Group Relative Policy Optimization}
  \textbf{Input} initial policy model $\pi_{\theta_{\text{init}}}$; reward models $r_\varphi$; task prompts $\mathcal{D}$; 
  hyperparameters $\varepsilon$, $\beta$, $\mu$
  \begin{algorithmic}[1]
    \State Initialize policy model: $\pi_\theta \leftarrow \pi_{\theta_{\text{init}}}$
    \For{iteration = 1, \dots, I}
       \State Set reference model: $\pi_{ref} \leftarrow \pi_{\theta}$
      \For{step = 1, \dots, M}
      \State Draw a mini-batch $\mathcal{D}_b$ from $\mathcal{D}$
      \State Assign old policy model: $\pi_{\theta_{old}} \leftarrow \pi_{\theta}$ 
      \State For every question $q \in \mathcal{D}_b$, generate $G$ outputs $\{o_i\}_{i=1}^G \sim \pi_{\theta_{old}} (\cdot \mid q)$
      \State Evaluate each generated output $o_i$ using $r_{\varphi}$ to obtain rewards $\{r_i\}_{i=1}^{G}$
      \State Estimate the group-relative advantages $\hat{A}_{i,t}$ for the $t$-th token of every $o_i$
      \For{GRPO iteration = 1, \dots, $\mu$}
        \State Update the policy model $\pi_{\theta}$ by optimizing the GRPO objective (Equation \ref{eq:GC-GRPO})
      \EndFor
    \EndFor
    \State Perform continual learning for $r_\varphi$ with a replay buffer. 
    \EndFor 
  \end{algorithmic}
  \textbf{Output} $\pi_\theta$
  \label{alg:iter-grpo}
\end{algorithm}

\subsubsection{Outcome Supervision RL with GRPO} More formally, for each prompt $q$, a set of candidate responses $\{o_1, o_2, \cdots, o_G\}$ are generated by sampling from $\pi_{\theta_{old}}$. Each response is evaluated by the reward model, resulting in $G$ reward values $\mathbf{r}=\{r_1, r_2, \cdots, r_G\}$. The obtained rewards are standardized by subtracting the batch mean and dividing by the batch standard deviation. With outcome supervision, the normalized reward is only assigned at the completion of each output $o_i$ and the same normalized value is used as the advantage for every token, i.e., $\hat{A}_{i, t} = \widetilde{r}_i = \frac{r_i- {\rm mean}(\mathbf{r})}{{\rm std}(\mathbf{r})}$ for all $t$. The policy parameters are then updated by maximizing the objective in equation (\ref{eq:GRPO-obj}).

\subsubsection{Process Supervision RL with GRPO} Outcome-level reward signals alone may not be sufficient for guiding policies in tasks requiring multi-step reasoning. Inspired by \cite{wang2023math}, process supervision is considered, assigning feedback after each individual reasoning step. More specifically, given prompt $q$ and sampled responses $\{o_1, o_2, \cdots, o_G\}$, a process reward model generates stepwise reward signals: $\mathbf{R} = \{ \{r_1^{index(1)},\cdots,r_1^{index(K_1)}\}, \cdots,  \{r_G^{index(1)},\cdots,r_G^{index(K_G)}\} \}$, where $index(j)$ denotes the end-of-step token index for the $j$-th step, and $K_i$ is the number of steps in response $o_i$. As before, the rewards are normalized with respect to their group, via $\widetilde{r}_i^{index(j)} = \frac{r_i^{index(j)} - {\rm mean}(\mathbf{R})}{{\rm std}(\mathbf{R})}}$. 
In this setting, the advantage attributed to each token is calculated as the sum of the normalized stepwise rewards for all steps yet to occur, i.e., $\hat{A}_{i, t} = \sum_{index(j) \ge t} \widetilde{r}_i^{index(j)}$. The policy optimization proceeds by maximizing the target in equation (\ref{eq:GRPO

\subsubsection{Iterative RL with GRPO} During reinforcement learning, the previously trained reward model may become inadequate for guiding the evolving policy model. To address this, we implement an iterative reinforcement learning scheme using GRPO. As detailed in Algorithm \ref{alg:iter-grpo}, this approach involves creating updated reward model training datasets from samples produced by the current policy model. The reward model is incrementally updated using a replay buffer that includes 10\% of prior training examples. Subsequently, the policy model is assigned as the new reference, and training continues for the policy model with supervision from the updated reward model.

\subsection{Training and Evaluating DeepSeekMath-RL}

Reinforcement learning is performed using DeepSeekMath-Instruct 7B as the base. The RL training set comprises chain-of-thought-style questions drawn from GSM8K and MATH, originating from the SFT data, totaling approximately 144K instances. To assess the influence of RL on datasets not exposed during training, other SFT examples are excluded. Reward model training data is constructed in accordance with the method outlined by \citep{wang2023math}. The initial reward model utilizes DeepSeekMath-Base 7B and is trained with a learning rate of 2e-5. For GRPO, the policy model is updated with a learning rate of 1e-6 and a KL penalty coefficient of 0.04. Each question prompts $64$ sampled completions, with a maximum sequence length of 1024 and a batch size of 1024 for training. After each round of exploration, the policy model is updated only once. DeepSeekMath-RL 7B is assessed on various benchmarks consistent with the evaluation protocol for DeepSeekMath-Instruct 7B. Within DeepSeekMath-RL 7B, tasks from GSM8K and MATH formatted with chain-of-thought reasoning are considered in-domain, while all other test sets are treated as out-of-domain.

Table \ref{tab:sft_rl_math} presents a comparative analysis of the performance of both open- and closed-source models equipped with chain-of-thought and tool-augmented reasoning strategies on benchmarks in English and Chinese. Key observations include: 1) When applying chain-of-thought reasoning, DeepSeekMath-RL 7B achieves accuracy rates of 88.2\% on GSM8K and 51.7\% on MATH. These results outperform all other open-source models ranging from 7B to 70B parameters, as well as most closed-source models. 2) Notably, the DeepSeekMath-RL 7B model is fine-tuned exclusively with chain-of-thought-format instruction data from GSM8K and MATH, initialized from DeepSeekMath-Instruct 7B. Despite the limited variety in its training corpus, DeepSeekMath-RL 7B consistently exceeds the performance of DeepSeekMath-Instruct 7B across all assessment criteria, highlighting the effectiveness of reinforcement learning.

\section{Discussion}
This section discusses insights derived from our pre-training and reinforcement learning (RL) experiments.

\subsection{Lessons Learnt in Pre-Training}

We begin by sharing reflections on our pre-training process. Unless stated otherwise, training procedures follow those detailed in Section \ref{sec:quality-policy}. For this analysis, the term DeepSeekMath Corpus refers to the version comprising 89 billion tokens from the second iteration of data collection.

\subsubsection{Code Training Benefits Mathematical Reasoning}

An often-cited but largely untested belief posits that training on code data enhances reasoning skills. Here, we seek to partially address this claim within the mathematical context, providing evidence that code pre-training boosts a model’s proficiency in mathematical reasoning tasks, both in standard and tool-assisted scenarios.

To systematically examine the impact of code-based pre-training on mathematical reasoning, we adopted both two-stage and one-stage training frameworks:

\textbf{Two-Stage Training}

\begin{itemize}[topsep=0pt]
    \item \textbf{Code Training for 400B Tokens $\rightarrow$ Math Training for 150B Tokens}: DeepSeek-LLM 1.3B is initially exposed to 400B code tokens, after which it undergoes an additional 150B tokens of math-focused training;
    \item \textbf{General Training for 400B Tokens $\rightarrow$ Math Training for 150B Tokens}: For comparison, we carry out a baseline experiment using general tokens—sourced from a broad, large-scale corpus compiled by DeepSeek-AI—instead of code tokens in the first training stage, to evaluate the specific benefits of code-related pretraining versus general data for improving mathematical reasoning skills.
\end{itemize}

\textbf{One-Stage Training}

\begin{itemize}[topsep=0pt]
    \item \textbf{Math Training for 150B Tokens}: DeepSeek-LLM 1.3B is trained directly on 150B math tokens in a single phase;
    \item \textbf{Training on a mixture of 400B Code Tokens and 150B Math Tokens}: Sequential math training after code pretraining has been observed to diminish code task performance. Therefore, we also explore a concurrent training scheme where code and math tokens are jointly sampled within a unified training stage, aiming to enhance both mathematical reasoning and code-related capabilities, as well as to address issues of catastrophic forgetting.
\end{itemize}

\paragraph{Results}
The comparative results from various training methodologies are presented in Table \ref{tab:code-math} and Table \ref{tab:code-math-general-eval}.

Pretraining with code has a positive impact on program-aided mathematical problem-solving, evident in both staged and combined training regimens. Table \ref{tab:code-math} shows that when code pretraining precedes math specialization in the two-stage setup, there is a marked improvement in performance on tasks like GSM8K and MATH using Python. A subsequent round of math-focused training yields additional gains. Notably, with mixed one-stage training, where code and math data are integrated, the model suffers less from the typical loss of coding ability encountered in sequential setups. Furthermore, this strategy enhances outcomes on both code-centric tasks (Table \ref{tab:code-math-general-eval}) and program-aided mathematical reasoning benchmarks (Table \ref{tab:code-math}).

Beyond tasks involving external tools, code pretraining still confers notable gains in mathematical reasoning. Specifically, in the two-phase scheme, an initial code training period moderately elevates mathematical abilities and primes the model for more effective subsequent math training, resulting in peak performance. However, a joint one-stage approach with both code and math data appears detrimental for tool-free mathematical reasoning, which may be attributable to the constrained capacity of DeepSeek-LLM 1.3B, potentially preventing it from effectively assimilating both domains at once.

\subsubsection{Limited Efficacy of ArXiv Papers in Enhancing Mathematical Reasoning}

Although arXiv papers frequently constitute a substantial fraction of math pre-training datasets \citep{minerva,gpt-f,llemma,mathpile}, there has been relatively little granular investigation into how effectively they improve mathematical reasoning skills. Somewhat surprisingly, our empirical findings indicate that exposure to arXiv papers does not appear to bolster mathematical reasoning capabilities. Our experiments encompass a range of model scales, such as DeepSeek-LLM 1.3B and DeepSeek-Coder-Base-v1.5 7B \citep{deepseek-coder}, utilizing variously curated arXiv text sources:

\begin{itemize}[topsep=0pt]
    \item \textbf{MathPile} \citep{mathpile}: a collection comprising 8.9B tokens filtered and cleaned via heuristic processes, with more than 85\% sourced from scientific arXiv manuscripts;
    \item \textbf{ArXiv-RedPajama} \citep{redpajama}: a dataset consisting of the complete set of arXiv LaTeX source files, preprocessed to exclude preambles, comments, macros, and bibliographies, encompassing 28.0B tokens in total.
\end{itemize}

We carried out experiments wherein DeepSeek-LLM 1.3B was trained for 150B tokens and DeepSeek-Coder-Base-v1.5 7B for 40B tokens, each on their respective arXiv-derived datasets. The results suggest that pre-training on arXiv papers fails to yield significant gains, and may even have detrimental effects, on a wide spectrum of mathematical reasoning benchmarks. This observation holds consistently across several task types considered, including quantitative reasoning (e.g., GSM8K and MATH as reported in Table \ref{tab:arxiv-cot}), multiple-choice assessments such as MMLU-STEM (Table \ref{tab:arxiv-cot}), and formal mathematics problems like those in miniF2F (Table \ref{tab:arxiv_atp}).

Nevertheless, it is important to acknowledge several caveats to this interpretation. The following areas remain to be systematically investigated:

\begin{itemize}[topsep=0pt]
    \item The role of arXiv-derived data in supporting niche or underexplored mathematical tasks, such as converting formal theorem statements or proofs into informal text (informalization), which falls outside the purview of this study;
    \item The potential synergistic effects when arXiv data is integrated with diverse pre-training data sources;
    \item Whether scaling models further could reveal advantages not apparent at the current sizes.
\end{itemize}

Given these open questions, a more comprehensive evaluation is warranted, which we propose to pursue in subsequent work.

\subsection{Reinforcement Learning Perspectives}
\subsubsection{Toward a Unified Framework} Here, we set forth a unified framework to systematically analyze a variety of training techniques—including SFT, RFT, DPO, PPO, and GRPO—and empirically probe the impact of several factors within this framework. In general, the gradient with respect to the parameters $\theta$ for a given training strategy can be expressed as follows:

\begin{equation}
    \nabla_{\theta}\mathcal{J}_{\textcolor{red}{\mathcal{A}}}(\theta) = \mathbb{E}[\underbrace{(q,o) \sim \textcolor{red}{\mathcal{D}}}_{Data \ Source}]\left( \frac{1}{|o|} \sum_{t=1}^{|o|}  \underbrace{GC_{{\mathcal{A}}}(q, o, t, \textcolor{red}{\pi_{{rf}}})}_{Gradient \ Coefficient}  \nabla_{\theta}\log \pi_{\theta}(o_t | q, o_{<t})\right).
\label{eq:objective}
\end{equation}

This framework consists of three fundamental elements: 1) \textit{Data Source $\mathcal{D}$}, which provides the samples for model training; 2) \textit{Reward Function $\pi_{{rf}}$}, responsible for supplying the reward signals that guide learning; and 3) \textit{Algorithm $\mathcal{A}$}, which integrates the training samples and reward feedback to compute the gradient coefficient $GC$, thereby regulating the intensity of reinforcement or penalization for each data point. We discuss a set of illustrative algorithms formulated within this general schema:

\begin{itemize}
    \item  \textbf{Supervised Fine-tuning (SFT)}: SFT further trains a pretrained model using a curated set of SFT data selected by humans.
    \item \textbf{Rejection Sampling Fine-tuning (RFT)}: Building on SFT, RFT continues training on outputs sampled from the SFT model in response to SFT prompts, but only includes samples passing a correctness-based filtering step.
    \item \textbf{Direct Preference Optimization (DPO)}: DPO advances SFT by updating the model with additional, pairwise-labeled outputs obtained from the SFT model, optimizing the model with a DPO loss that reflects preferences.
    \item \textbf{Online Rejection Sampling Fine-tuning (Online RFT)}: Unlike standard RFT, Online RFT begins with the SFT model and performs continual refinement, repeatedly sampling and fine-tuning on newly generated outputs from the evolving policy model itself.
    \item \textbf{PPO/GRPO}: PPO/GRPO also use the SFT model as initialization, but employ reinforcement signals to further improve the model based on outputs dynamically generated by the policy model in real time.
\end{itemize}

A summary of these method components is presented in Table \ref{tab:data_policy}. For a comprehensive explanation of the derivation, consult Appendix \ref{app:analysis-rl}.

\paragraph{Discussion on Data Source} We classify data sources into two main types: online sampling and offline sampling. Online sampling refers to collecting training data from the ongoing exploration of the current policy model during training, whereas offline sampling utilizes data obtained from the initial SFT model’s outputs. RFT and DPO utilize an offline sampling framework, whereas Online RFT and GRPO implement an online approach.

Figure \ref{fig:rl-analysis} demonstrates that Online RFT achieves substantially higher performance than RFT on two evaluation benchmarks. Notably, both methods yield similar results in the initial training phase; however, as training progresses, Online RFT achieves a clear performance edge, illustrating the benefits of online sampling. This phenomenon can be explained by the fact that, at the outset, the actor closely mirrors the SFT model, leading to only slight discrepancies in their sampled data. As training continues, data sampled from the evolving actor displays increasingly pronounced differences, amplifying the benefits provided by continuously updated, real-time data.

\paragraph{Discussion on Gradient Coefficient} The algorithms translate reward signals into gradient coefficients, which in turn inform model parameter updates. In our experiments, the reward function is either a ‘Rule’ or a ‘Model’: the former assesses response quality according to answer correctness, while the latter employs a trained reward model to assign scores to responses, using training data that is itself labeled via the rule-based method. Equations \ref{eq:GC-RFT} and \ref{eq:GC-GRPO} illuminate a fundamental distinction: GRPO dynamically modifies its gradient coefficient in line with reward model outputs, enabling finer-grained reinforcement or penalization across different response qualities. Conversely, Online RFT does not incorporate this differentiation—incorrect responses are not penalized, and every correct response receives identical reinforcement strength.

As illustrated in Figure \ref{fig:rl-analysis}, GRPO consistently outperforms online RFT, underscoring the effectiveness of adjusting both positive and negative gradient coefficients. Moreover, GRPO+PS yields better outcomes than GRPO+OS, which suggests that leveraging more granular, step-sensitive gradient coefficients offers clear advantages. We further examine the effect of iterative RL by performing two cycles of training in our experiments. According to Figure \ref{fig:iter-rl}, iterative RL leads to substantial performance gains, particularly during the initial iteration.

\subsubsection{Why RL Works?} This study employs reinforcement learning on a subset of instruction tuning data, resulting in notable improvements over the instruction tuning model baseline. To deepen our understanding of RL's effectiveness, we assess both Pass@K and Maj@K accuracy for the Instruct and RL models across two benchmarks. As presented in Figure \ref{fig:maj-pass}, RL considerably boosts Maj@K while leaving Pass@K largely unaffected. These results suggest that reinforcement learning strengthens the output distribution, specifically by amplifying the likelihood that a correct response is found among the top-K outputs, rather than advancing the model's core competencies. Likewise, \citep{wang2023making} observed a \textbf{misalignment problem} in reasoning within SFT models, demonstrating that their performance can be augmented using a range of preference alignment techniques \citep{yuan2023rrhf,song2023preference,wang2023making}.

\subsubsection{How to Achieve More Effective RL?}
\label{sec:effec_rl}
We show that RL delivers strong performance gains on mathematical reasoning benchmarks, and we present a comprehensive framework unifying various notable training approaches. Within this framework, every method can be understood as a form of direct or approximate RL. Equation \ref{eq:objective} outlines three principal elements: Data Source, Algorithm, and Reward Function, which may guide future research.

\paragraph{Data Source} The choice of data source underpins every training approach. For RL in particular, the data source refers to unlabeled queries and their sampled outputs generated by the policy model. In the current study, we restrict ourselves to questions originating from the instruction tuning phase and utilize a basic nucleus sampling method. This limitation may explain why the RL pipeline leads to improved Maj@K performance exclusively. Looking ahead, we intend to apply the RL pipeline to novel, out-of-distribution prompts and incorporate \textbf{advanced sampling (decoding) strategies}, such as tree-search algorithms \citep{yao2023tree}. In addition, adopting \textbf{efficient inference techniques} \citep{xia-etal-2023-speculative,leviathan2023fast,kwon2023efficient,xia2024unlocking}—which impact the exploration proficiency of the policy model—will likely be crucial.

\paragraph{Algorithms} Algorithms utilize the data and reward signal to compute gradient coefficients, which in turn adjust the model parameters. Referring to Equation \ref{eq:objective}, current approaches largely place complete \textbf{TRUST} in the reward function’s signals, leveraging them to either boost or diminish the conditional probability assigned to specific tokens. Nonetheless, there is no guarantee that these reward signals remain accurate at all times, particularly when handling highly complex tasks. For instance, despite the rigorous curation by expert annotators, the PRM800K datasets \citep{lightman2023let} are reported to still possess around 20\% erroneous annotations\footnote{\url{https://github.com/openai/prm800k/issues/12\#issuecomment-1728491852}}. This highlights the necessity of investigating reinforcement learning algorithms that demonstrate robustness in the face of noisy reward signals. We posit that alignment strategies grounded in \textbf{WEAK-TO-STRONG} principles \citep{burns2023weak} may yield transformative improvements to foundational learning algorithms.

\paragraph{Reward Function} The reward function acts as the primary provider of learning signals during training. Within the RL framework, this function is typically embodied as a neural reward model. We identify three major avenues for advancing reward models: 1) \textbf{Improving the reward model's generalization capability.} Ensuring the reward model effectively transfers to novel, out-of-distribution prompts and sophisticated decoding outputs is crucial; otherwise, reinforcement learning may do little more than maintain the status quo in LLM behavior, rather than drive substantive advancements; 2) \textbf{Capturing the uncertainty inherent in the reward model.} Quantifying and leveraging uncertainty could serve as an important interface between imperfect reward models and weak-to-strong learning frameworks; 3) \textbf{Developing efficient, high-fidelity process reward models} that deliver precise, fine-grained feedback to guide the model’s reasoning steps \citep{lightman2023let,wang2023math}.

\section{Conclusion, Limitation, and Future Work} 

This work introduces DeepSeekMath, a model that surpasses all other open-source approaches on the competition-grade MATH benchmark and nearly matches the performance of proprietary models. DeepSeekMath~is built upon DeepSeek-Coder-v1.5 7B and is continually trained on 500B tokens, notably including 120B math-focused tokens extracted from Common Crawl. Through comprehensive ablation analysis, we demonstrate that web-based data sources offer considerable promise for obtaining high-quality mathematical datasets, whereas arXiv did not prove as advantageous as initially anticipated. We propose Group Relative Policy Optimization (GRPO)—an adaptation of Proximal Policy Optimization (PPO)—which enhances mathematical reasoning ability while consuming less memory. Our experimental findings indicate that GRPO remains beneficial for DeepSeekMath-Instruct 7B, even at elevated benchmark scores. Furthermore, we establish a unified framework for interpreting several methodologies and highlight multiple promising avenues for advancing reinforcement learning strategies.

Despite DeepSeekMath~achieving strong results on quantitative reasoning tasks, it currently lags behind closed-source models when tackling geometry and theorem-proving. Specifically, in preliminary assessments, DeepSeekMath~struggled with geometry problems involving triangles and ellipses—likely reflecting a bias in the data used during pre-training and fine-tuning. Additionally, limited by its model size, DeepSeekMath~does not match GPT-4’s capacity for few-shot learning; GPT-4 exhibits clear gains with few-shot inputs, while DeepSeekMath~shows no significant difference between zero-shot and few-shot scenarios. Looking ahead, we plan to further enhance our curated data selection system to assemble a richer and more robust pre-training dataset. Additionally, we intend to explore further research directions (Section \ref{sec:effec_rl}) aimed at realizing more efficient reinforcement learning paradigms for large language models (LLMs).

\end{document}